\documentclass[aspectratio=169, 14pt]{beamer}
\input{preamble}

\title{From function to agent}

\begin{document}

% --- Slide 1: Section divider ---
\sectiondivider{From function to agent}

% --- Slide 2: Web vs terminal ---
\begin{frame}{Two ways in}
  \vfill
  \centering
  \begin{tikzpicture}[node distance=4cm]
    \node[tjo box fill, minimum width=4.5cm, minimum height=2cm] (web)
      {\textbf{Web interface}\\[4pt]{\normalsize ChatGPT, Claude.ai}\\[2pt]
       {\small\color{tjo@black!60} chat UI hides the API}};
    \node[tjo box fill, minimum width=4.5cm, minimum height=2cm, right=3cm of web] (term)
      {\textbf{Terminal}\\[4pt]{\normalsize cURL, Python}\\[2pt]
       {\small\color{tjo@black!60} direct API access}};

    % Arrow
    \draw[tjo arrow, <->] (web.east) -- (term.west)
      node[midway, above, font=\normalsize] {same API underneath};
  \end{tikzpicture}
  \vfill
\end{frame}

% --- Slide 3: The agent loop (code) ---
\begin{frame}[fragile]{The agent loop}
  \vfill
\begin{codeblock}
while not done:
    response = llm(system_prompt + history)
    tool_calls = parse(response)
    results = execute(tool_calls)
    history += [response, results]
\end{codeblock}

  \vskip1.5em
  {\color{tjo@darkteal}\bfseries
    This is the complete architecture of every coding agent}
  \vfill
\end{frame}

% --- Slide 4: The agent loop (diagram) ---
\begin{frame}{The agent loop}
  \vfill
  \centering
  \begin{tikzpicture}[node distance=1.6cm, every node/.style={align=center}]
    % System prompt + history at top
    \node[tjo box, fill=tjo@darkteal, text=white, minimum width=4cm,
          minimum height=1cm, font=\normalsize]
      (sys) {System prompt + History};

    % LLM
    \node[tjo box fill, below=1.2cm of sys, minimum width=4cm,
          minimum height=1cm, font=\normalsize]
      (llm) {\textbf{LLM}\; {\small\code{f(s) $\to$ s}}};

    % Parse
    \node[tjo box, right=3cm of llm, minimum height=1cm, font=\normalsize]
      (parse) {Parse tool calls};

    % Execute
    \node[tjo box, below=1.2cm of parse, minimum height=1cm, font=\normalsize]
      (exec) {Execute tools};

    % Arrows
    \draw[tjo arrow] (sys) -- (llm);
    \draw[tjo arrow] (llm) -- (parse)
      node[midway, above, font=\small] {response};
    \draw[tjo arrow] (parse) -- (exec)
      node[midway, right, font=\small] {call};

    % Results arrow: loop back around into LLM
    \draw[tjo arrow] (exec.south) -- ++(0,-0.6) -| ([xshift=-0.5cm]llm.south)
      node[near start, below, font=\small] {results};

    % Feedback to history
    \draw[tjo arrow, dashed, color=tjo@cyan]
      ([xshift=0.5cm]llm.north) -- ([xshift=0.5cm]sys.south)
      node[midway, right, font=\small, text=tjo@cyan] {append};
  \end{tikzpicture}
  \vfill
\end{frame}

% --- Slide 5: Tools ---
\begin{frame}{Tools}
  \vfill
  \centering
  You describe tools in the system prompt

  \vskip1.5em
  The LLM generates structured text requesting a tool call

  \vskip1.5em
  Your scaffolding parses and executes it
  \vfill
\end{frame}

% --- Slide 6: The key insight ---
\statementslide{The LLM never executes anything}

% --- Slide 7: Two tools ---
\begin{frame}[fragile]{Two tools}
  \vfill
  \centering
\begin{codeblock}
read_file(path)           -> contents
write_file(path, content) -> ok
\end{codeblock}

  \vskip1.5em
  {\color{tjo@darkteal}\bfseries That is enough to do useful work}
  \vfill
\end{frame}

% --- Slide 8: Production equivalence ---
\statementslide{Claude Code, Cursor, Copilot:\\[0.8em]
  the same loop, more tools,\\[0.3em] better prompts}

% --- Slide 9: Live demo ---
\statementslide{{\fontsize{36}{44}\selectfont LIVE DEMO}\\[1em]
  {\fontsize{20}{26}\selectfont\mdseries Watch the agent work}}

\end{document}
